\bigskip
\bigskip
\subsubsection*{Խնդիրներ և հարցեր թեմա 7-ի վերաբերյալ}

\begin{enumerate}[label=\thesection.\arabic*.]

% 14.1
\item Ճի՞շտ է արդյոք, որ կամայական $X$ և $Y$ տարածությունների դեպքում $P_X: X \times Y \to X$ կանոնական պրոյեկցիան փակ արտապատկերում է:

\begin{hint}
Տես օրինակ 3-ը թեմա 13-ում:
\end{hint}

% 14.2
\item Ապացուցեք. Կամայական $X$ և $Y$ տարածությունների դեպքում $X \times Y$-ը հոմեոմորֆ է $Y \times X$-ին:

\begin{hint}
Հիմնավորեք, որ $f: X \times Y \to Y \times X, f(x,y)=(y,x)$ արտապատկերումը հոմեոմորֆիզմ է:
\end{hint}

% 14.3
\item Ապացուցեք. Կամայական $X, Y, Z$ տարածությունների դեպքում $(X \times Y) \times Z$-ը հոմեոմորֆ է $X \times (Y \times Z)$-ին:

\begin{hint}
Հիմնավորեք, որ $g: (X \times Y) \times Z \to X \times (Y \times Z), g((x,y),z)=(x,(y,z))$ արտապատկերումը հոմեոմորֆիզմ է:
\end{hint}

% 14.4
\item Գոյություն ունի՞ $(X, \tau)$ տարածություն, որը հոմեոմորֆ է իր $(X \times X, \tau \times \tau)$ քառակուսուն:

\begin{hint}
Հիմնավորելով, որ $X$-ը պետք է լինի անվերջ բազմություն, դիտարկեք, օրինակ, բնական թվերի բազմությունը դիսկրետ տոպոլոգիայով:
\end{hint}

% 14.5
\item Ճի՞շտ է արդյոք, որ կամայական $X$ անվերջ բազմություն դիսկրետ տոպոլոգիայով հոմեոմորֆ է իր $(X \times X, \text{դիսկր.} \times \text{դիսկր.})$ քառակուսուն:

% 14.6
\item Դիտարկենք $A=[-1,1] \times [-1,1]$ քառակուսին $\mathbb{R}^2$ հարթությունում և նրան ներգծված $B=\{(x_1, x_2); x_1^2+x_2^2 \le 1\}$ շրջանը: Ապացուցեք, որ դրանք հոմեոմորֆ տարածություններ են:

\begin{hint}
Նախ համոզվեք, որ $A$-ի եզրագծի ամեն մի $(x_1, x_2)$ կետի դեպքում կամ $|x_1|=1, |x_2| \le 1$, կամ $|x_2|=1, |x_1| \le 1$: Հիմնվելով դրա վրա ապացուցեք. $h: A \to B, h(x_1, x_2) = \frac{\max\{|x_1|, |x_2|\}}{\sqrt{x_1^2+x_2^2}}(x_1, x_2)$ արտապատկերումը հոմեոմորֆիզմ է:
\end{hint}

% 14.7
\item Ապացուցեք. $\mathbb{R}^n$ էվկլիդեսյան մետրիկական տարածության $\mathbb{R}^n \setminus \{0\}$ ենթատարածությունը հոմեոմորֆ է $S^{n-1} \times \mathbb{R}$ արտադրյալին, որտեղ $0$-ն $\mathbb{R}^n$-ի սկզբնակետն է, իսկ $S^{n-1}$-ը՝ $\{x=(x_1, \dots, x_n); \|x\| = \sqrt{\sum x_i^2}=1\}$ սֆերան է:

\begin{hint}
Հիմնավորեք, որ $f: \mathbb{R}^n \setminus \{0\} \to S^{n-1} \times \mathbb{R}, f(x)=(\frac{1}{\|x\|}(x_1, \dots, x_n), \lg \|x\|)$ արտապատկերումը հոմեոմորֆիզմ է:
\end{hint}

% 14.8
\item Կամայական $f: X \to Y$ արտապատկերման համար սահմանվում է նրա գրաֆիկը որպես $X \times Y$ արտադրյալի $\Gamma_f = \{(x,y); y=f(x)\}$ ենթաբազմություն: Ապացուցեք.
\begin{enumerate}
    \item[ա)] եթե $f$-ը անընդհատ է, ապա $\Gamma_f$-ը որպես $X \times Y$-ի ենթատարածություն հոմեոմորֆ է $X$-ին:
    \item[բ)] եթե $Y$-ը հաուսդորֆյան տարածություն է, ապա $\Gamma_f$-ը փակ ենթաբազմություն է $X \times Y$-ում:
\end{enumerate}

% 14.9
\item Ապացուցեք. $X$-ը հաուսդորֆյան տարածություն է այն և միայն այն դեպքում, երբ $X \times X$-ի $\Delta = \{(x,x); x \in X\}$ անկյունագիծը փակ ենթաբազմություն է $X \times X$-ում:

\begin{hint}
Օգտվել նախորդ խնդրի բ) պնդումից որպես $f$ վերցնելով $X$-ի նույնական արտապատկերումը:
\end{hint}

\end{enumerate}

\par Մինչև հաջորդ խնդիրներին անցնելը տանք մի սահմանում: Ասում են որ տոպոլոգիական տարածություններին առնչվող հատկությունը պահպանվում է ուղիղ արտադրյալի դեպքում, եթե ամեն անգամ, երբ որևէ երկու տարածություններ բավարարում են այդ հատկությանը, նրանց ուղիղ արտադրյալը նույնպես բավարարում է այդ հատկությանը:
\par Օրինակ, թեորեմ 7-ից հետևում է, որ անջատելիության $T_2$ աքսիոմը պահպանվում է ուղիղ արտադրյալի դեպքում: Մյուս կողմից որպես գիտելիք նշենք, որ նորմալությունը (որպես անջատելիության աքսիոմ) չի պահպանվում ուղիղ արտադրյալի դեպքում:

\begin{enumerate}[label=\thesection.\arabic*., resume]

% 14.10
\item Ապացուցեք, որ $T_0, T_1$ աքսիոմները պահպանվում են ուղիղ արտադրյալի դեպքում:

% 14.11
\item Ապացուցեք, որ սեպարաբելությունը պահպանվում է ուղիղ արտադրյալի դեպքում:

% 14.12
\item Դիցուք $B_1=\{U_i; i \in I\}$ և $B_2=\{V_j; j \in J\}$ ընտանիքները տոպոլոգիայի բազիս են համապատասխանաբար $(X, \tau)$ և $(Y, \mathcal{T})$ տարածությունների համար: Ապացուցեք. $B=\{U_i \times V_j; i \in I, j \in J\}$ ընտանիքը $(X \times Y, \tau \times \mathcal{T})$ տարածության տոպոլոգիայի բազիս է:

% 14.13
\item Ապացուցեք. հաշվելիության երկրորդ աքսիոմը պահպանվում է ուղիղ արտադրյալի դեպքում:

% 14.14
\item Ապացուցեք. հաշվելիության առաջին աքսիոմը պահպանվում է ուղիղ արտադրյալի դեպքում:

\begin{hint}
Նախ ձևակերպեք և ապացուցեք 14.12 խնդրին նմանակ պնդում կետի շրջակայքների բազիսի վերաբերյալ:
\end{hint}

\end{enumerate}
